\section{From the matrix path to the action}\label{sec:action-limit}

We now prove Theorem~\ref{thm:action}. The matrix comparison controls points
where $M(B)=aa^{\transpose}$ exactly. The differential equation supplies a
penalty for the remaining points, and compactness makes that penalty
effective without requiring convergence of the chosen coefficients $C$.

\begin{proof}[Proof of Theorem~\ref{thm:action}]
Fix $\delta>0$, and use the field $M$ and primitive $J$ from
Section~\ref{sec:field}. For each $0<\varepsilon\le1$, the connecting-path
result of Section~\ref{sec:path} provides a lower triangular matrix
$C_\varepsilon$ and a path $B_\varepsilon$ satisfying
\eqref{eq:framework-path}. Write $a_\varepsilon=C_\varepsilon u$ and temporarily
suppress the subscripts. The path $B$ is $C^1$, and $a\in H^1_0(I;\R^N)$
is absolutely continuous. The chain and product rules give, almost everywhere,
\begin{align}
 (J(B)+a^{\transpose}Ba)'
 &=2(a')^{\transpose}Ba
      -\tr\bigl((M(B)-aa^{\transpose})B'\bigr)\notag\\
 &=2(a')^{\transpose}Ba
      -\varepsilon^{-1}\|M(B)-aa^{\transpose}\|_{\HS}^2.
 \label{eq:action-identity}
\end{align}
Since $a$ vanishes at the endpoints, integration and
$2(a')^{\transpose}Ba\le |a'|^2+a^{\transpose}B^2a$ yield
\begin{equation}\label{eq:penalty-bound}
 J(-K)-J(K)\le\mathcal E_K(a)
      +\int_I F_\varepsilon(B(x),a(x))\,dx,
\end{equation}
where
\[
 F_\varepsilon(B,a)
 =|a|^{2q}-a^{\transpose}(K^2-B^2)a
       -\varepsilon^{-1}\|M(B)-aa^{\transpose}\|_{\HS}^2.
\]

The bounds from Section~\ref{sec:path} place every pair
$(B_\varepsilon(x),a_\varepsilon(x))$, for $x\in\overline I$ and
$0<\varepsilon\le1$, in a single compact set $\mathcal K$.
This uses the continuous representative of $u$, the bound on $C_\varepsilon$,
and fixed $\delta$. We claim that
\begin{equation}\label{eq:uniform-penalty}
 \eta_\varepsilon:=\max_{(B,a)\in\mathcal K}
       \max\{F_\varepsilon(B,a),0\}\longrightarrow0.
\end{equation}
Otherwise there would be $\eta>0$, a sequence $\varepsilon_n\downarrow0$,
and a convergent sequence $(B_n,a_n)\in\mathcal K$ with
$F_{\varepsilon_n}(B_n,a_n)\ge\eta$.
The first two terms of $F_\varepsilon$ are uniformly bounded on
$\mathcal K$, so $M(B_n)-a_na_n^{\transpose}\to0$.
At the limit $(B_*,a_*)$ we have $M(B_*)=a_*a_*^{\transpose}$ and
\[
 |a_*|^{2q}-a_*^{\transpose}(K^2-B_*^2)a_*\ge\eta,
\]
contrary to Proposition~\ref{prop:rankone}. This proves
\eqref{eq:uniform-penalty}.

Equation~\eqref{eq:penalty-bound} now implies
\[
 J(-K)-J(K)\le
 \sup_{C\text{ lower triangular}}\mathcal E_K(Cu)+|I|\eta_\varepsilon.
\]
Let $\varepsilon\downarrow0$ at this fixed $\delta$. By
Lemma~\ref{lem:jensen},
\[
 \sup_C\mathcal E_K(Cu)
 \ge\sum_i\int_{-k_i}^{k_i}
       \bigl((k_i^2-s^2+\delta)^\sigma-\delta^\sigma\bigr)\,ds.
\]
Finally let $\delta\downarrow0$. On each integration interval the integrand
converges to $(k_i^2-s^2)^\sigma$ and lies between zero and this same function,
by subadditivity of the power $\sigma$. Dominated convergence gives
\eqref{eq:action-bound}. No convergence of $C_\varepsilon$, and no bound
uniform in $\delta$, has been used.
\end{proof}
